% =====================================================================
\chapter{Modelos de espacio de estados}
\label{tema:5}
% =====================================================================

La atención guarda todo el pasado y lo consulta en cada paso; una red recurrente lo resume en un estado de tamaño fijo. Los modelos de espacio de estados (SSM) son redes recurrentes lineales con una estructura heredada de la teoría de sistemas: se definen en tiempo continuo, se discretizan y pueden calcularse como recurrencia, con coste constante por token en la generación, o como convolución, en paralelo durante el entrenamiento. Mamba los hace selectivos, haciendo que sus parámetros dependan de la entrada, y recupera el cálculo paralelo mediante un barrido asociativo. Se sigue a Gu, Goel y Ré~\cite[§2]{s4}, a Gu y Dao~\cite[§2 y §3]{mamba}, la parte de la clase~10 de CS336 dedicada a las alternativas al Transformer~\cite{cs336} y la sesión de A.~Gu en el seminario CS25~\cite{cs25}. Los sistemas discretos y su respuesta al impulso siguen el curso 6.003 del MIT~\cite{mit6003}.

\section{Sistemas lineales continuos y su discretización}

\begin{definicion}[Modelo de espacio de estados]\label{def:ssm}
Sean $A\in\R^{n\times n}$, $B\in\R^{n\times 1}$ y $C\in\R^{1\times n}$. El \emph{modelo de espacio de estados} continuo con entrada escalar $x(t)$, estado $h(t)\in\R^{n}$ y salida $y(t)$ es
\[
h'(t)=Ah(t)+Bx(t),\qquad y(t)=Ch(t)
\]
\cite[ecs.~1a y~1b]{mamba}. Un modelo de lenguaje aplica un SSM independiente a cada una de las $d$ componentes de la secuencia de vectores.
\end{definicion}

Para procesar una secuencia se fija un paso $\Delta>0$ y se supone que la entrada es constante en cada intervalo de longitud $\Delta$ (\emph{retención de orden cero}).

\begin{proposicion}[Discretización por retención de orden cero]\label{prop:zoh}
Sea $n=1$, con $A=a\ne 0$ y $B=\beta$. Si $x(\tau)=x_t$ es constante en $[t\Delta,(t+1)\Delta)$, entonces $h_{t+1}=\Abar h_t+\Bbar x_t$, con
\[
\Abar=e^{a\Delta},\qquad \Bbar=\frac{e^{a\Delta}-1}{a}\,\beta .
\]
Si $a<0$, entonces $0<\Abar<1$. En el caso matricial, $\Abar=\exp(\Delta A)$ y $\Bbar=(\Delta A)^{-1}\bigl(\exp(\Delta A)-I\bigr)\Delta B$~\cite[ec.~4]{mamba}.
\end{proposicion}

S4 usa en su lugar la discretización bilineal, $\Abar=(I-\frac\Delta2A)^{-1}(I+\frac\Delta2A)$ y $\Bbar=(I-\frac\Delta2A)^{-1}\Delta B$~\cite[§2.3]{s4}.

\addtocounter{definicion}{1}%

\section{Recurrencia y convolución}

\begin{definicion}[SSM discreto]\label{def:ssmdiscreto}
El SSM discreto con matrices $(\Abar,\Bbar,C)$ y estado inicial $h_0=0$ transforma la sucesión $(x_t)_{t\ge1}$ en
\[
h_t=\Abar h_{t-1}+\Bbar x_t,\qquad y_t=Ch_t
\]
\cite[ecs.~2a y~2b]{mamba}.
\end{definicion}

\begin{proposicion}[Forma convolucional]\label{prop:conv}
Si $(\Abar,\Bbar,C)$ no dependen de $t$, la salida del SSM discreto es la convolución causal
\[
y_t=\sum_{j=0}^{t-1}\bar K_j\,x_{t-j},\qquad \bar K_j=C\Abar^{\,j}\Bbar,
\]
y $\bar K$ es la respuesta del sistema al impulso unitario ($x_1=1$ y $x_t=0$ para $t>1$).
\end{proposicion}

\addtocounter{definicion}{1}%

Las dos formas tienen costes complementarios. La recurrencia cuesta $O(n)$ operaciones por token y un estado de tamaño fijo, como una red recurrente, pero es secuencial. La convolución se calcula en paralelo sobre toda la secuencia, por ejemplo con la transformada rápida de Fourier en $O(T\log T)$, que es como S4 entrena~\cite[§2.4]{s4}. Gu y Dao llaman a esta propiedad \emph{invariancia temporal lineal}~\cite[§2]{mamba}.

\section{Selección: el modelo de Mamba}

La invariancia temporal impide que el modelo decida, según el contenido, qué recordar y qué olvidar: el núcleo $\bar K$ trata igual cada token. Gu y Dao argumentan que esta es la razón por la que los SSM anteriores eran débiles en lenguaje, y proponen que los parámetros dependan de la entrada~\cite[§3.1 y §3.2]{mamba}.

\begin{definicion}[SSM selectivo]\label{def:selectivo}
Un \emph{SSM selectivo} usa en el instante $t$ las matrices $B_t=s_B(x_t)$, $C_t=s_C(x_t)$ y el paso $\Delta_t=\softplus\bigl(p+s_\Delta(x_t)\bigr)$, donde $s_B$, $s_C$ y $s_\Delta$ son proyecciones lineales de la entrada y $p$ es un parámetro, y discretiza con la \cref{prop:zoh} en cada instante~\cite[§3.2, algoritmo~2]{mamba}. La matriz $A$ es fija y diagonal.
\end{definicion}

\begin{teorema}[Selección y puertas]\label{teo:puerta}
Sea un SSM selectivo escalar con $A=-1$, $B=1$ y $\Delta_t=\softplus(z_t)$, donde $z_t$ es una función lineal de $x_t$. Su recurrencia es
\[
h_t=(1-g_t)\,h_{t-1}+g_t\,x_t,\qquad g_t=\sigma(z_t),
\]
la de una red recurrente con una puerta que interpola entre conservar el estado y sustituirlo por la entrada~\cite[§3.5.1, teorema~1]{mamba}.
\end{teorema}

\addtocounter{definicion}{1}%

\section{El barrido paralelo}

Con parámetros dependientes de $t$, el coeficiente de $x_{t-j}$ en $h_t$ es $\Abar_t\cdots\Abar_{t-j+1}\Bbar_{t-j}$, que depende de $t$: ya no hay núcleo fijo ni convolución. La linealidad en el estado basta, sin embargo, para paralelizar.

\begin{proposicion}[Asociatividad de la composición de pasos]\label{prop:scan}
Sea $\phi_t(h)=a_th+b_t$ el paso $t$ de una recurrencia escalar, representado por el par $(a_t,b_t)$. Se define
\[
(a,b)\bullet(a',b')=(a'a,\ a'b+b').
\]
Entonces $(\phi_{t'}\circ\phi_t)$ está representado por $(a_t,b_t)\bullet(a_{t'},b_{t'})$, la operación $\bullet$ es asociativa y, con $h_0=0$, $h_t$ es la segunda componente de $(a_1,b_1)\bullet\cdots\bullet(a_t,b_t)$.
\end{proposicion}

Por la \cref{prop:scan}, los estados $h_1,\dots,h_T$ son los prefijos de un producto asociativo, que un algoritmo de barrido paralelo calcula con profundidad $O(\log T)$ y trabajo $O(T)$. Mamba combina este barrido con la fusión de las operaciones en la memoria rápida de la GPU y el recálculo de los estados intermedios en la pasada hacia atrás, sin escribir nunca el estado expandido en la memoria principal~\cite[§3.3.2]{mamba}. Con estado vectorial y $A$ diagonal, cada componente es una recurrencia escalar independiente.

\section{Comparación con la atención}

\begin{table}[ht]
\centering\small
\caption{Atención causal frente a un SSM selectivo, por capa, con $T$ tokens y dimensión $d$.}\label{tab:ssm}
\begin{tabular}{@{}lll@{}}
\toprule
& Atención & SSM selectivo \\
\midrule
Coste del entrenamiento & $O(T^{2}d)$ & $O(Td)$ por cada dimensión de estado \\
Memoria en la generación & caché que crece con $T$ & estado de tamaño fijo \\
Coste por token generado & $O(Td)$ & independiente de $T$ \\
Recuperar un dato exacto lejano & directo & limitado por la compresión del estado \\
\bottomrule
\end{tabular}
\end{table}

Gu y Dao informan de un rendimiento de inferencia unas cinco veces mayor que el de un Transformer de tamaño similar~\cite[§4.5]{mamba}. La clase~10 de CS336 señala la debilidad complementaria: los SSM anteriores a Mamba resolvían mal las tareas de recuperación asociativa, importantes en lenguaje, y los diseños recientes son híbridos que conservan algunas capas de atención; Jamba, por ejemplo, intercala capas Transformer y Mamba en proporción $1:7$ dentro de una MoE de $5{,}2\cdot 10^{10}$ parámetros~\cite{cs336,jamba}. No hay información pública de que Claude, Gemini o Grok usen SSM.

